\documentclass[12pt,a4paper]{article}
\usepackage {amsmath,amsthm}
\usepackage{geometry}
\usepackage{titlesec}
 \geometry{a4paper, total={170mm,257mm}, left=20mm, right=20mm, top=25mm, bottom=25mm}
\usepackage{sectsty}
\usepackage{natbib}
\usepackage[hidelinks]{hyperref}
\usepackage{multirow}
\usepackage{makecell}
\usepackage{adjustbox}
%\usepackage{ulem}
\usepackage{threeparttable}
\usepackage{longtable}
\usepackage{soul}
\usepackage[title]{appendix}
%\usepackage{prettyref}

%\usepackage[colorlinks=true, urlcolor=Blue]{hyperref}
%\usepackage{titling}
\sectionfont{\fontsize{14}{15}\selectfont}
\subsectionfont{\fontsize{13}{15}\selectfont}

\titleformat*{\section}{\LARGE\bfseries}
\titleformat*{\subsection}{\Large\bfseries}


\usepackage{caption}
\captionsetup{belowskip=0pt,aboveskip=6pt}

\makeatletter
\setlength{\@fptop}{0pt}
\makeatother

\makeatletter
\newcommand*\exInput[1]{\@@input#1 }
\makeatother


\usepackage[dvipsnames]{xcolor}
\newcommand{\fnote}[1]{\captionsetup{font={small, it}, aboveskip=0pt} \caption*{#1}}
\newcommand{\link}[1]{{\color{blue}\href{#1}{#1}}}



%\usepackage{cite}
\usepackage{booktabs}
\usepackage{float}
\usepackage{setspace}
\usepackage{placeins}
\usepackage[list=true]{subcaption}
\captionsetup[sub]{font=footnotesize}

\usepackage{graphicx}
\usepackage{fvextra}


\newtheorem{theorem}{Proposition}
\newtheorem{hypothesis}{Hypothesis}
	\def\hypothesisautorefname{Hypothesis}
\newtheorem{result}{Result}
	\def\resultautorefname{Result}
\def\sym#1{\ifmmode^{#1}\else\(^{#1}\)\fi}


%\title{\Large Willingness to Pay for Signals of Rare Events\\}
\title{\Large Consistency Training\\}

\author{Alexander Ugarov\thanks{Email: \href{mailto:augarov@hivereview.org}{augarov@hivereview.org}.} \\
{\small \textit{Hivereview}}
}

\begin{document}
\maketitle
\onehalfspacing
\begin{abstract}
\noindent 
Reinforcement Learning from internal feedback (RLIF) promises potential gains in areas where human feedback is scarce but typically delivers only fleeting marginal gains at the cost of entropy collapse. We improve upon this method by calculating internal rewards from a more easily verifiable part of the completion and using a more accurate label-free judge. We train a label-free judge to distinguish between solutions to reasoning problems produced from intact vs corrupt reasoning traces. The judge recovers more than 80\% of the accuracy gap between random pick and verifiable rewards in best-of-8 selection for the Countdown task. Our early results show that rejection fine-tuning using internal feedback indeed increases accuracy on Countdown task by 11 pp. for Llama-3.2-7B model and by 16 pp. for Qwen3-0.6B which approaches the gains from verifiable rewards.


\vspace{6pt}
\noindent {\footnotesize Keywords: reinforcement learning, GRPO, self-verification, LLM-as-a-judge, reward hacking, label-free rewards}
\end{abstract}
\normalsize

\section{Introduction}
\label{sec:intro}
Reinforcement learning of large language models mostly utilizes verifiable rewards and human-generated data. As models approach superhuman abilities in selected domains, designing reinforcement learning environments and human verification in general becomes challenging. This problem raises the interest for methods which can improve models’ performance without verifiable data even if potential improvements are modest.

Human learning is not necessarily limited by verifiable rewards. People can often improve skills just by practicing them. For example, just doing homework or test problems without feedback can improve the performance on subsequent tests (\citet{fyfe_mathematics_2017}, \citet{avvisati_learning_2020}). There are at least two reasons for why practicing is helpful on its own. First, practicing often involves reading and learning more. Second, students can often understand whether they manage to solve a problem. When solving a challenging problem, people can experience an Aha! moment in which a sudden insight delivers a solution with high degree of confidence. Hence the process of solution as well as process of finding mistakes in own reasoning constitutes an additional source of feedback.

We explore whether reasoning models can also learn from "Aha!" moments even if rewards are not immediately verifiable. Our method falls into the stream of literature called Reinforcement Learning From Internal Feedback (RLIF), the term first introduced by \citet{zhao_learning_2026}. \citet{agarwal_unreasonable_2025} find that RL with rewards using only the answer's entropy significantly increase Qwen-7B accuracy on multiple math and coding benchmarks.  \citet{kang_scalable_2025} propose an alternative measure of internal feedback (self-certainty) which is based on average cross-entropy of policy with respect to the uniform distribution and improves best-of-N selection for diverse reasoning benchmarks. \citet{zhao_learning_2026} use self-certainty reward in GRPO to train models in math and coding. 

However, many of these results are suspect due to using early base Qwen models only which can improve with spurious rewards \citep{shao_spurious_2026}. Additionally, \citet{zhang_no_2025} find that Instruct models have much smaller improvements from common RLIF method and this improvement turns negative when the policy inevitably start gaming the rewards. They also show that using either self-consistency or trajectory entropy as rewards leads to entropy collapse due to positive and high covariance between advantages and original policy.


Our approach negates this channel of entropy collapse common in RLIF by splitting completion into a reasoning trace and a final solution and calculating rewards only on the solution part. Hence rewards do not push reasoning to conform to the existing policy and reasoning trace constitutes the largest portion of completions in our setting. Additionally, we use different policies for generation and for scoring solutions.


%The completions we explore include reasoning trace, final clean solution, and the answer. Getting clean step-by-step solutions is the immediate target of the training while effective reasoning trace is the main goal. The training process rewards reasoning traces for their ability to end up in a solution which is considered to be highly likely by a model. This setup most fits the domains where typical reasoning involves extensive backtracking and parallel exploration, while the solution is short and easy to verify (puzzles, mathematics, software engineering). 

%We close generation-verification gap through a careful design of verification metric which uses no human-labelled data. Similar to \citet{agarwal_unreasonable_2025}, we use the total log-prob of a completion as a proxy for the output’s correctness. However, we specifically train a judge model to align log-probabilities with probabilities of an answer being correct using  

We find that models are more accurate in verifying solutions rather than full completions containing reasoning traces, solutions , and answers. However, moderate-sized LLM tend to be poor judges and require additional training in order to enable self-play training. The paper explores two potential solutions. First, fine-tuning a model with the judge prompt and gold answers marginally improves its ability to discriminate between correct and incorrect solutions. Second, training a model to recognize a solution produced from corrupted reasoning produces a more sizable gain in verification abilities.

Training with self-plausibility and format rewards using a label-free judge increases accuracy on Countdown task from  46\% to 53\% on  Llama-3.2-3B-Instruct model and from 22\% to 32\% for Qwen3-0.6B model. This improvement beats training on format-rewards only which increases the accuracy only by 0.5 pp (decreased by 10 pp. for the Qwen3 model) with a difference significant at 5\%. The increase is comparable to potential gains with verifiable rewards and happens without significant decrease in response diversity. 

Our two contributions are separating the more easily verifiable part of the completion (cleaned solution) from the traces where RLIF training can have more negative effect on solution diversity, and the design of label-free judge with more accurate scoring of solutions.

\vspace{10pt}
\section{Preliminaries}

\textbf{Task} Each Countdown tasks involves finding a sequence of arithmetic operations such that it produces the expected number using a given set of numbers such that each number is used only once. Our tasks mostly come from TinyZero dataset (citation) and include tasks involving just 3 numbers and tasks involving 4 numbers which are more complex. Many tasks have multiple solutions, and the number of solutions can be used as a proxy of their difficulty.

The advantage of Countdown task for our purpose is that it is much easier to verify solution than to find it. It makes reasoning for solving the problem which is similar to more complex mathematics and coding tasks. At the same time, even small models can manage to verify solutions.


% Trajectory-level entropy
\textbf{Reward.} Our \textit{plausibility reward} is the total log-probability assigned by a judge model to the part of the completion which includes solution and the final answer but omits the reasoning:

\begin{equation}
R_p=\sum_{k=k_0}^S \log(p_J(x_k|x_{<k}))
\end{equation}

Here $k_0$ represents the first token of the generated solution, $p_J$ is the judge policy, and $p_J(x_k|x_{<k})$ is the judge model's probability of token $x_k$ conditional on the history.

When $J$ is the same generation/sampling policy, and the span covers the whole completion, the metric equivalent to \textit{trajectory-level entropy} is introduced by \citet{agarwal_unreasonable_2025}. For a pretrained model the trajectory-level entropy of a completion converges to the probability of observing this completion in the data given the prompt. However, this calibration generally breaks during post-training (\citet{openai_gpt-4_2024}).

We train the judge using  Direct Preference Optimization  (DPO) (\citet{rafailov_direct_2024}) to align generated log-probabilities with log-probabilities of completions being correct. Suppose $x$ is a prompt coming from distribution $D$, and $y_w$ and $y_l$ are two different completions with $y_w$ being the preferred solution. DPO minimizes the following objective (negative log-likelihood) with respect to policy $\pi(\cdot | x)$: 

\begin{equation}
\min_{\pi(\cdot |x)} \mathcal{L}=\min_{\pi(\cdot |x)} E[-\log \sigma(r(x,y_w)-r(x,y_l))]
\end{equation}

Here the $\sigma()$ is a logistic function $\sigma(x) \equiv {1\over 1+e^{-x}}$ and the reward function is:
\begin{equation}
r(x,y)=\beta \log {\pi(y|x)\over \pi_{ref}(y|x)}, \beta>0
\end{equation}

In our judge training, preferred solutions $y_w$ in each pair come from full reasoning traces while rejected solutions $y_l$ come from corrupt reasoning traces. Corruption involves either empty traces or significantly cut traces. In this setup, the (log) trained judge's policy is a linear combination of the base (log) policy with  with log-likelihood ratio of winning vs losing solution.

\begin{theorem} Let preferred solution $y_w \sim p_w(y|x)$ and the losing solution $y_l \sim p_l(y|x)$. Assume that these distributions have the same support and $y_w$ and $y_l$ are independent given $x$. Then the following  policy minimizes DPO objective in the class of all possible policies (for any $y$ in support of $p_l(y|x),p_w(y|x)$):
\begin{equation}
\log \pi^*(y|x) = \log \pi_{ref}(y|x)+{1\over \beta} \log \left({p_w(y|x)\over p_l(y|x)}\right) +Z(x)
\end{equation}
\end{theorem}

\begin{proof}
Consider drawing a pair $(y_1,y_2)$ such that both $(y_1,y_2)=(y_w,y_l)$ and $(y_1,y_2)=(y_l,y_w)$ are equally likely events. Let $T\sim\operatorname{Bernoulli}(1/2)$ be an indicator variable such that \((Y_1,Y_2)=(Y_w,Y_l)\) if \(T=1\) and \((Y_1,Y_2)=(Y_l,Y_w)\) otherwise, and assume that $T$ is independent of $y_w,y_l$ draws. By Bayes rule, the conditional probability of the event $T=1$ is:

$$P(T=1|y_1,y_2)={p_w(y_1|x)p_l(y_2|x) \over p_w(y_1|x)p_l(y_2|x) +p_l(y_1|x)p_w(y_2|x) }=\sigma(\rho(y_1)-\rho(y_2))$$

Here $\rho(y)\equiv \log {p_w(y|x)\over p_l(y|x)}$.  Define $s(y_1,y_2)\equiv\sigma(r(x,y_1)-r(x,y_2))$ (omitting $x$ for brevity). It follows from the definition that $s(y_2,y_1)=1-s(y_1,y_2)$. The original DPO objective can be written as:
$$\mathcal{L}=E_{y_w,y_l}[-\log s(y_w,y_l)]=E_{T,y_1,y_2}[-T\log s(y_1,y_2)-(1-T)\log(1-s(y_1,y_2))]$$

Applying the law of total expectation:
$$\mathcal{L}=E_{y_1,y_2}[-P(T=1|y_1,y_2)\log s(y_1,y_2)-P(T=0|y_1,y_2)\log(1-s(y_1,y_2))]$$

Note that $0<s(y_1,y_2)<1$ and hence the expression in square brackets is cross-entropy of distribution $s()$ with respect to the conditional probabilities $P(T=1|y_1,y_2), P(T=0|y_1,y_2)$. Hence by Gibbs inequality $s(y_1,y_2)=P(T=1|y_1,y_2)$ minimizes $\mathcal{L}$. It implies:

$$\sigma(r(x,y_1)-r(x,y_2))=\sigma(\rho(y_1)-\rho(y_2))$$

Given that $\sigma(\cdot)$ is a strictly monotonic function, we need that for each $x$, $y_1$ and $y_2$:
$$r(x,y_1)-r(x,y_2)=\rho(y_1)-\rho(y_2)$$

Hence any optimal solution $r^*(x,y)$ has to satisfy the following condition:
\begin{equation}
r^*(x,y)=\rho(y)+c(x)=\log \left({p_w(y|x)\over p_l(y|x)}\right)+c(x)
\end{equation}

Here $c(x)$ is a function of $x$ as any function of $x$ wouldn't affect the objective (but it still need to satisfy normalization constrain given that $\pi^*(y|x)$ is a probability distribution). Hence every optimal DPO policy satisfies:
\begin{equation}
\log \pi^*(y|x) = \log \pi_{ref}(y|x)+{1\over \beta} \log \left({p_w(y|x)\over p_l(y|x)}\right) +{1\over \beta} c(x)
\end{equation}

Renaming $Z(x)={1\over \beta} c(x)$ obtains the equation in the statement.
\end{proof}


The assumption of expressiveness of policy then guarantees that there is a policy $\pi^*()$ corresponding to the reward function $r^*(x,y)$. Note, that the log-likelihood ratio is a noisy signal of solution accuracy:
\begin{equation}
\log {p_w(y|x)\over p_l(y|x)} = \log {P_w(z(y)|x)\over P_l(z(y)|x)}+\log{p_w(y|x, z(y))\over p_l(y|x, z(y))}
\end{equation}

Here $z(y)$ is a truth value of completion $y$ and $P_w(z|x), P_l(z|x)$ are the probabilities to obtain this truth value given $x$. The second component of this decomposition measures the stylistic contribution while the first one measures the actual correctness. If there are no stylistic differences between solutions coming from corrupt and non-corrupt traces, the stylistic component becomes zero, and the ratio measures log-odds of $p_w$ being a correct solution.

\vspace{10pt}

\section{Setup}
\label{sec:setup}

We select two models for our tests: Llama-3.2-3B-Instruct  and Qwen3-0.6B. The Llama model is an instruct-type model, which means it doesn't adhere to any reasoning format which would conflict with our experiments.  However, our tests with the Llama model show very low initial accuracy on Countdown with a zero-shot prompt and no reasoning effort. The Qwen3-0.6B model combines better and more modern training with small size. Larger Qwen3 models saturate Countdown tasks for 4 digit problems.


For each model, we prepare a generating model and a judge model as two adapters of the same base model. The judge adapter comes from Direct Preference Optimization (DPO) where each training pair consists of one preferred solution block coming from an intact reasoning trace and another - from corrupt reasoning trace. Generating model's adapter comes from supervised fine-tuning on synthetic reasoning traces which include backtracking and verification behaviors. Generating model adapter is necessarily only because the base model has poor adherence to format (Qwen3) or no reasoning behavior (Llama-3.2-3B-Instruct). This stage could be potentially replaced by training on gold reasoning completions, completions from a stronger model, or dropped completely if the original model is strong enough for RL stage.


\textbf{Generating Model Training.}  Following \citet{gandhi_cognitive_2025} we seed reasoning behavior for each base model by fine-tuning it on reasoning traces containing backtracking, verification of results, and repeatedly taking stock of numbers already used while solving the puzzle. Procedural generation builds all the 6000 reasoning traces for fine-tuning to guarantee good representation for all the desired reasoning behaviors. These examples use the same prompt as the one used in our plausibility training which asks for a reasoning part, solution, and the answer. Fine-tuning also covers 3000 gold answers with a neutral system (no format restrictions) to prepare the model for judge prompts. All the puzzles come from TinyZero dataset (\citet{pan_tinyzero_2025}) with equal split to 3-number and 4-number puzzles. 

The training runs for one epoch using effective batch size of 16 samples. The learning rate follows cosine schedule with 5\% warmup and peaks at  1e-4. The increase in Countdown performance is significant. On 3-number puzzles the greedy accuracy goes from 8\% to 72\% for the Llama model and from 2\% to 38\% for the Qwen3 model. On 4-number puzzles the accuracy increases from 0\% to 24\% for the Llama and 0\% to 12\% for the Qwen3  (both numbers use 50 puzzles per size).\footnote{Evaluation use the same prompt consisting of the puzzle followed by format instructions, without a worked example. Examples however do not improve scoring for the untrained model.} The improvement happens both due to better format following and better reasoning.



\textbf{Judge training}.  The model for scoring solutions (judge) is trained with Direct Preference Optimization (DPO) (\citet{rafailov_direct_2024}) using solution and answer parts of completions  produced by the fine-tuned generation model. Preferred solutions are the original solutions generated with full reasoning traces.  Rejected solutions come from corrupted reasoning traces. We use 1800 reasoning completions produced by a fine-tuned model to create two pairs for DPO training per completion. The first pair's corrupt solution comes from empty reasoning trace. The second pair's corrupt solution comes from only the first half of the reasoning trace.  Hence rejected solutions use much less reasoning to get to the solution and tend to be less accurate.

The training starts with the fine-tuned generation model and continues for one epoch with the effective batch of 16 pairs. It uses learning rate of 2e-5 with cosine schedule and $\beta=0.1$.



\textbf{Reinforcement Learning with Plausibility Rewards.} We train the generating model for 150 steps using Group Relative Preference Optimization algorithm (GRPO) with both format and plausibility rewards. Format rewards incentivize completions with non-verbal answers and three tagged parts (reasoning, solution, answer). Formatting rewards are essential for our setup because the calculation of rewards requires having a formatted solution while accuracy verification needs numeric-only format for our grader.\footnote{The grader's scores do not affect the training.}. \textit{Plausibility reward} is the sum of logits of the solution+answer parts of the completion as graded by the judge model using a neutral system prompt (see "Preliminaries" section above). We calibrate the weight of plausibility rewards to match the standard deviation of formatting rewards.

We use TRL library for training,  and vLLM library for generation and plausibility scoring. All the training runs use Low Rank Adaptaton (LoRA) \citep{hu_lora_2021} with rank of 16 and $\alpha=32$ using PEFT library. We observe no significant improvement from increasing LORA rank to 32. The training uses learning rate of 2e-05 with constant schedule after 20 warm-up steps.\footnote{Some charts show 300 steps for the plausibility rewards with the last 150 steps coming from an extra run starting from step 150 checkpoint.}  GRPO uses 16 completions per step with anchoring to the base model ($\beta_{KL}=0.02$) and no rewards scaling to avoid large changes in gradient when diversity of responses becomes small.




\vspace{10pt}

\section{Results}
\label{sec:trained-judges}
\textbf{Scoring Solutions.} Even without a trained judge, scoring solutions and answers identifies correct solutions more reliably than scoring whole completions (Table 1). For example, SFT policy for the Llama model picks a correct solution with probability 0.803 if seeing the whole completion vs 0.838 for the solution block only. The base model identifies correct solution in the pair with probability 0.768 when using the whole completion\footnote{This is not surprising as it doesn't expect reasoning.}, but with probability 0.863 when using solution and answer. This evidence supports our decision to score only the solution part of completions.

% Generated by data/scope_countdown/make_judges_countdown.py (do not edit by hand: rerun it) -> judges_countdown.csv.
% Score = sum (the GRPO reward's reduction); the mean-score version: tables/archive/tab_judges_countdown_mean.tex.
% Llama: context_scores.csv (2,560 GRPO training completions of run mean-34 + cont100b on the mixed SFT; the
% completions of tab_context_countdown). Qwen: Drive countdown/sft/Qwen3-0.6B-...-imp0x0-mix3x500/
% judge_corrupt_empty-cut0.1-cut0.25-cut0.5_unfinished_n4x1000_n3x1000_t0.9_cap2048_seed0/dpo_empty_b0.1_lr2e-05_ep2/
% calibration_sum_missmed_S_grade_n200x8_t0.9_cap1024_seed0_scores.csv (copied to data/scope_countdown/qwen/).
\begin{table}[H]
\centering
\caption{Judges on Countdown: how often the \texttt{sum} score ranks correct completions of a puzzle above wrong ones.}
\label{tab:judges-countdown}
\begin{threeparttable}
\small
\begin{tabular}{llcccc}
\toprule
 & & \multicolumn{2}{c}{Llama-3.2-3B} & \multicolumn{2}{c}{Qwen3-0.6B} \\
\cmidrule(lr){3-4} \cmidrule(lr){5-6}
Judge & Scored span & Pairs & \makecell{Groups\\ separated} & Pairs & \makecell{Groups\\ separated} \\
\midrule
Base model policy & whole completion & 0.768 & 0.222 & 0.515 & 0.170 \\
Base model policy & solution + answer & 0.863 & 0.593 & 0.841 & 0.610 \\
SFT policy & whole completion & 0.803 & 0.284 & 0.610 & 0.229 \\
SFT policy & solution + answer & 0.838 & 0.481 & 0.799 & 0.525 \\
Judge policy & solution + answer & 0.965 & 0.790 & 0.960 & 0.834 \\
\midrule
Shorter solution & solution + answer & 0.651 & 0.148 & 0.782 & 0.453 \\
\bottomrule
\end{tabular}
\begin{tablenotes}\footnotesize
\item Score: \texttt{sum}, the total token log-probability of the scored span. Solution + answer is scored after the judge prompt (neutral system prompt and the bare puzzle), the whole completion after the policy's prompt (puzzle and format instructions). Pairs: share of same-puzzle (correct, wrong) pairs ranked correctly (ties count 1/2). Groups separated: share of puzzles with both correct and wrong completions in which every correct completion scores above every wrong one.
\end{tablenotes}
\end{threeparttable}
\end{table}


The judge scoring is significantly more accurate in identifying correct solutions. Within a pair, it picks the correct solution with probability of 0.965 (0.863 for the base model policy). It allows to completely separate correct vs incorrect solutions in almost 80\% of groups of 8 sampled completions. The judge performance is comparable between the two models. Hence it can provide informative training signals for the next stage of our experiment - GRPO training. 


\label{sec:corrupt}
\textbf{Accuracy Improvement.}  GRPO training significantly increases the sampled accuracy on Countdown both for Llama and Qwen3 models. The starting (SFT) model's accuracy with pass@1 at temperature 0.9 is just 0.459 for the Llama and 0.224 for Qwen3-0.6B models using the equal mix of 3-number and 4-number puzzles. Training with plausibility and format rewards for 150 steps improves the accuracy to 0.528 for Llama and 0.322 for the Qwen model.

Given evidence of GRPO improving accuracy with spurious rewards, we train the model using format-only rewards as a control. With Llama model, format-only training doesn't affect the accuracy, while it actually hurts the accuracy of Qwen3 model. Hence learning format cannot explain the observed improvements from GRPO training with plausibility rewards.

We observe small reduction in solution diversity for both models throughout the training. Llama model experiences pass@8 decrease from 0.58 to 0.51 for 4-number puzzles and no decrease for 3-number puzzles. Training Qwen3 reduces pass@8 metric from 0.88 to 0.75 for 3-number and from 0.32 to 0.275 for 4-number puzzles. The number of distinct correct solutions also falls down by roughly 14-20\% for both models.

% Source: Drive countdown/grpo/Llama-3.2-3B-Instruct-countdown3000-ep1-lr0.0001-seed0-tr2-judge3000-mix3x3000-judge_dpo-start_sft/runs/
%   control:   grpo-format_only-mix3_0.5-expr-kl0.02-format-noscale-mixv3
%   treatment: grpo-pass-fp0.05-mix3_0.5-expr-kl0.02-corrupt_empty-cut-v3-mix-corrupt-pass
% Both: mixed SFT (3,000 + 3,000 version-2 traces + 3,000 judge-format examples), TRAIN_MIX_3 0.5, KL 0.02, 16 x 8 per
% step, cosine LR, warmup 20, 150 steps, SCALE_REWARDS False, expression reward on. Treatment adds the corrupted-reasoning
% judge dpo_empty-cut (DPO, chosen = solution from normal reasoning, rejected = from empty / cut reasoning; no checker
% labels), pass arm: reward 1 if the judge's mean log-prob of solution + answer is above the threshold at which 5% of the
% starting policy's wrong calibration answers pass. Data: ../data/grpo_format_vs_corrupt/ (summary.csv; its test3 / test4mix / test5 columns are not used: in both runs
% those evals ran after [checkpoint eval] had loaded checkpoint 50, so they measure step 50); start evals in
% ../data/grpo_format_vs_solver/start/. Paired statistics: narrative.md, "format_only without SD scaling" (2026-10-27).
% Greedy row (2026-10-28): the 200 test puzzles (eval_test_*) + the 200 held-out puzzles of the held-out curve
% (eval_heldout_heldout_* + eval_heldout_unscored_*, never trained on), final weights; start: eval_test_before_n200_max1024
% + heldout_before_grpo_split_n4000_nopool_seed1/. Paired bootstrap over puzzles (10,000 draws, seed 0), exact sign test.
\begin{table}[H]
\centering
\caption{GRPO on Countdown (Llama-3.2-3B, 150 steps): format-only control vs the label-free corrupted-reasoning judge.}
\label{tab:format-vs-corrupt}
\begin{threeparttable}
\begin{tabular}{lcccc}
\toprule
 & Start & Format only & \makecell{Format +\\ label-free judge} & Difference \\
\midrule
\multicolumn{5}{l}{\emph{Sampled, $T=0.9$}} \\
pass@1 & 0.243 & 0.249 & 0.385 & 0.136 {\footnotesize [0.102, 0.171]} \\
pass@16 & 0.75 & 0.76 & 0.76 & 0.00 {\footnotesize [$-$0.08, 0.08]} \\
Distinct correct answers per puzzle & 2.83 & 3.01 & 2.70 & $-$0.31 {\footnotesize [$-$0.64, 0.01]} \\
\midrule
\multicolumn{5}{l}{\emph{Greedy}} \\
Accuracy & 0.265 & 0.248 & 0.315 & 0.068 {\footnotesize [0.030, 0.105]} \\
\bottomrule
\end{tabular}
\begin{tablenotes}\footnotesize
\item pass@$k$ and distinct correct answers: 16 samples at temperature 0.9 for each of 100 puzzles held out from the GRPO training draw. Greedy accuracy: 400 unseen 4-number puzzles (the 200 test puzzles and 200 held-out puzzles, including those 100). After 150 steps. Difference: judge run minus format only, paired over puzzles; 95\% confidence interval in brackets, from a bootstrap over puzzles (10,000 resamples; a puzzle's 16 samples stay together). pass@1 at step 100: difference 0.141 [0.106, 0.176]. Greedy, puzzles solved only by the judge run / only by the control: 45 / 18 (sign test $p = 0.001$). No reward scaling by the group standard deviation in either arm. One seed per arm.
\end{tablenotes}
\end{threeparttable}
\end{table}


We also train the model using rejection fine-tuning (RFT) approach \citep{yuan_scaling_2023}. Within each round, we sample 16 completions for 300 Countdown problems. For each incomplete solution without the solution block, we generate solution block separately after appending the incomplete completion to the reasoning prompt. For each group of 16 completions we select 2 completions with solutions having the highest plausibility reward. Within this group we drop 20\% of completions with the lowest score as they typically correspond to problems which the model doesn't manage to solve.\footnote{The accuracy drops only slightly without filtering bottom 20\% of solutions.} Then we fine-tune the model using our reasoning prompt and hyperparameters matching fine-tuning of the generating model described above (Section~\ref{sec:setup}). Training for 4 rounds exactly matches the amount of data seen through 150 steps of GRPO training.

RFT training results in higher accuracy gains for both models with less decrease in diversity. Llama model reaches sampled accuracy of 29\% on 4-number puzzles and 56\% overall. Trained Qwen3 model improves 4-number accuracy to 14\% and average accuracy to 38\%. Most of the learning happens before step 4, though Llama model continues to improve until step 7 which exceeds training budget of GRPO. Answer diversity as measured by pass@8 metric doesn't change until step 4 after which it starts to decrease. In contrast to GRPO, rejection fine-tuning also doesn't reduce the completion length and most of the completions do not fit into our budget of 1024 tokens. 

% Generated by data/rft_countdown/make_tab_rft.py (do not edit by hand: rerun it) -> per_puzzle.csv, summary.csv.
% Llama-3.2-3B RFT run: countdown/grpo/<Llama mixed SFT>-judge_dpo-start_sft/runs/
%   rft-judge-sum-missmed-corrupt_empty-cut-r8x300k2-skip0.2-lr2e-05-mix3_0.5-rft1-cut20 (judge sum, missing = median wrong)
% Qwen3-0.6B RFT run: countdown/grpo/qwen3-0.6b-sft-rft-judge-w3-p5-cut20 (judge sum, missing = 5th percentile)
% Both: 300 puzzles x 16 per round, top 2 finished above the group mean, bottom 20% of picks dropped, 1 epoch at 2e-5.
\begin{table}[H]
\centering
\caption{Rejection-sampling fine-tuning on Countdown with the label-free judge as the filter: pass@1 and pass@8.}
\label{tab:rft-countdown}
\begin{threeparttable}
\setlength{\tabcolsep}{4pt}
\begin{tabular}{lllcccc}
\toprule
Model & Metric & Puzzles & Start & Round 4 & Round 8 & Round 8 $-$ start \\
\midrule
\multirow{6}{*}{Llama-3.2-3B}
 & pass@1 & 3-number & 0.708 & 0.814 & 0.836 & 0.129 [0.095, 0.163] \\
 &  & 4-number & 0.210 & 0.293 & 0.291 & 0.081 [0.054, 0.110] \\
 &  & Both & 0.459 & 0.553 & 0.564 & 0.105 [0.083, 0.127] \\
\cmidrule(l){2-7}
 & pass@8 & 3-number & 0.980 & 0.975 & 0.960 & $-$0.020 [$-$0.050, 0.010] \\
 &  & 4-number & 0.580 & 0.575 & 0.570 & $-$0.010 [$-$0.075, 0.055] \\
 &  & Both & 0.780 & 0.775 & 0.765 & $-$0.015 [$-$0.050, 0.020] \\
\midrule
\multirow{6}{*}{Qwen3-0.6B}
 & pass@1 & 3-number & 0.368 & 0.641 & 0.622 & 0.254 [0.211, 0.299] \\
 &  & 4-number & 0.081 & 0.154 & 0.140 & 0.059 [0.033, 0.087] \\
 &  & Both & 0.224 & 0.398 & 0.381 & 0.157 [0.128, 0.184] \\
\cmidrule(l){2-7}
 & pass@8 & 3-number & 0.880 & 0.875 & 0.835 & $-$0.045 [$-$0.095, 0.005] \\
 &  & 4-number & 0.320 & 0.385 & 0.325 & 0.005 [$-$0.050, 0.055] \\
 &  & Both & 0.600 & 0.630 & 0.580 & $-$0.020 [$-$0.055, 0.015] \\
\bottomrule
\end{tabular}
\begin{tablenotes}\footnotesize
\item 200 test puzzles per size $\times$ 8 samples at temperature 0.9 (the puzzles of Table~\ref{tab:format-vs-corrupt}); start = the SFT model. Each round samples 300 new training puzzles $\times$ 16 completions, keeps per puzzle the two highest-scored finished completions above the group mean (dropping the lowest-scored 20\% of the round's picks), and fine-tunes on them for one epoch. Round 4 uses as many rollouts as 150 GRPO steps. Round 8 $-$ start: paired over puzzles, 95\% confidence interval in brackets.
\end{tablenotes}
\end{threeparttable}
\end{table}


\section{Future Work}
\label{sec:future}


\newpage

\bibliographystyle{plainnat-etal}
\bibliography{Consistency}

\begin{appendices}
\section{Tables}
\label{app:tables}
% Generated by data/grpo_countdown_150/make_tab_greedy_failures.py (do not edit by hand: rerun it) -> greedy_failures.csv.
% Llama-3.2-3B, the runs of tab_format_vs_corrupt (countdown/grpo/<Llama mixed SFT>-judge_dpo-start_sft/runs/):
%   format only grpo-format_only-mix3_0.5-expr-kl0.02-34-fmtw1-istt2-il-const;
%   judge grpo-sum-missmed-mix3_0.5-expr-kl0.02-corrupt_empty-cut-34-fmtw1-istt2-il-const. Start: eval_mixed_before_n200_max1024.
% Previous version (older runs, with the solver-pair judge): tables/archive/tab_greedy_failures_solver_pair.tex.
\begin{table}[H]
\centering
\caption{Llama-3.2-3B, greedy decoding after 150 GRPO steps: accuracy, truncation and made-up numbers.}
\label{tab:greedy-failures}
\begin{threeparttable}
\setlength{\tabcolsep}{4pt}
\begin{tabular}{lcccccc}
\toprule
 & \multicolumn{3}{c}{3-number puzzles} & \multicolumn{3}{c}{4-number puzzles} \\
\cmidrule(lr){2-4}\cmidrule(lr){5-7}
 & Start & \makecell{Format\\ only} & \makecell{Format +\\ label-free judge} & Start & \makecell{Format\\ only} & \makecell{Format +\\ label-free judge} \\
\midrule
Accuracy & 0.745 & 0.730 & 0.790 & 0.210 & 0.205 & 0.280 \\
\quad among finished & 0.805 & 0.730 & 0.790 & 0.300 & 0.206 & 0.280 \\
Truncated & 0.075 & 0.000 & 0.000 & 0.300 & 0.005 & 0.000 \\
Made-up numbers & 0.054 & 0.030 & 0.010 & 0.152 & 0.121 & 0.025 \\
\bottomrule
\end{tabular}
\begin{tablenotes}\footnotesize
\item The runs of Table~\ref{tab:format-vs-corrupt}, 200 test puzzles per size, one greedy completion each; start = the SFT model. Truncated: no answer within 1,024 tokens. Made-up numbers: share of answers using a number not among the given ones.
\end{tablenotes}
\end{threeparttable}
\end{table}

% Generated by data/synthetic_countdown/make_synthetic_table.py (do not edit by hand: rerun it). Source: results.csv
% from Consistency_Synthetic_Gaming_Countdown (2026-11-04); narrative.md, "Synthetic gaming tests on Countdown".
\begin{table}[H]
\centering
\caption{Synthetic gaming tests on Countdown (Llama-3.2-3B): share of puzzles in which the correct solution scores above its edited version.}
\label{tab:synthetic-countdown}
\setlength{\tabcolsep}{3pt}
\begin{adjustbox}{max width=\textwidth}
\begin{tabular}{lccccccccccccccc}
\toprule
 & \multicolumn{5}{c}{Base model} & \multicolumn{5}{c}{SFT policy} & \multicolumn{5}{c}{Corrupt DPO} \\
\cmidrule(lr){2-6}\cmidrule(lr){7-11}\cmidrule(lr){12-16}
Edit & \texttt{low10} & \texttt{mean} & \texttt{min} & SC & \texttt{sum} & \texttt{low10} & \texttt{mean} & \texttt{min} & SC & \texttt{sum} & \texttt{low10} & \texttt{mean} & \texttt{min} & SC & \texttt{sum} \\
\midrule
Other puzzle's solution & 1.00 & 1.00 & 0.69 & 1.00 & 1.00 & 1.00 & 1.00 & 1.00 & 1.00 & 1.00 & 1.00 & 1.00 & 0.99 & 1.00 & 1.00 \\
False last step & 0.86 & 0.85 & 0.50 & 0.81 & 0.87 & 0.90 & 0.91 & 0.81 & 0.96 & 0.90 & 0.98 & 0.98 & 0.95 & 0.97 & 0.98 \\
Unavailable number & 0.91 & 0.91 & 0.56 & 0.94 & 0.91 & 0.99 & 0.99 & 0.98 & 0.99 & 0.99 & 0.95 & 0.95 & 0.84 & 0.95 & 0.95 \\
Answer only & 0.03 & 1.00 & 0.01 & 0.58 & 0.01 & 1.00 & 1.00 & 1.00 & 1.00 & 1.00 & 1.00 & 1.00 & 0.97 & 1.00 & 1.00 \\
First step only & 0.94 & 1.00 & 0.49 & 0.99 & 0.87 & 1.00 & 1.00 & 1.00 & 1.00 & 1.00 & 1.00 & 1.00 & 0.97 & 1.00 & 1.00 \\
False last step + praise & 1.00 & 0.99 & 0.50 & 0.95 & 1.00 & 1.00 & 1.00 & 1.00 & 1.00 & 1.00 & 1.00 & 1.00 & 0.95 & 0.99 & 1.00 \\
Prose restatement & 1.00 & 0.60 & 0.00 & 0.00 & 1.00 & 1.00 & 1.00 & 1.00 & 1.00 & 1.00 & 1.00 & 1.00 & 0.90 & 1.00 & 1.00 \\
Redundant true step & 1.00 & 0.96 & 0.44 & 1.00 & 1.00 & 1.00 & 1.00 & 1.00 & 1.00 & 1.00 & 1.00 & 1.00 & 0.96 & 1.00 & 1.00 \\
\bottomrule
\end{tabular}
\end{adjustbox}

\smallskip
\begin{minipage}{\textwidth}\footnotesize
334 test puzzles (200 with 4 numbers, 134 with 3). Each correct solution is a solver-written route (one step per line, then ``Answer: \ldots''), checked by the grader, and is scored against each edit of it under the judge prompt (neutral system prompt and the bare puzzle; the scored text is the solution, the answer and the end-of-turn token). Edits: another puzzle's correct solution; a route over the given numbers that misses the target, with the last step's result replaced by the target; a correct-looking route that uses an unavailable number; the answer line alone; the first step and the answer; the false last step followed by ``All steps check out.''; the sentence ``We need to reach $T$ using $a$, $b$, $c$ and $d$.'' before the steps; the true step ``$T * 1 = T$'' after them. The last four edits keep a correct answer. Ties count 1/2. \texttt{mean} / \texttt{sum}: mean / total token log-probability; \texttt{low10}: sum of the 10 lowest token log-probabilities; \texttt{min}: the lowest token log-probability; SC: self-certainty, the mean over scored positions of the KL divergence from the uniform distribution to the model's next-token distribution. Base model: Llama-3.2-3B-Instruct; SFT policy: GRPO's starting point, fine-tuned on solver-written judge-format solutions in the same style; corrupt DPO: the corrupted-reasoning judge.
\end{minipage}
\end{table}

% Generated by data/grpo_countdown_150/make_tab_format_vs_corrupt.py (do not edit by hand: rerun it) -> per_puzzle.csv,
% summary.csv. 150 steps, constant LR 2e-5 after 20 warmup steps, KL 0.02, T 0.9; 200 + 200 test puzzles x 8, cap 1,024.
% Llama-3.2-3B runs (countdown/grpo/<Llama mixed SFT>-judge_dpo-start_sft/runs/): format only
%   grpo-format_only-mix3_0.5-expr-kl0.02-34-fmtw1-istt2-il-const; judge
%   grpo-sum-missmed-mix3_0.5-expr-kl0.02-corrupt_empty-cut-34-fmtw1-istt2-il-const ([test pass@k], checkpoint-150).
% Qwen3-0.6B runs: qwen3-0.6b-sft-format-n34-steps150-cap1024-istt2-seed0, qwen3-0.6b-sft-judge-w1-median (step 150).
\begin{table}[H]
\centering
\caption{GRPO on Countdown, 150 steps: pass@8 and distinct correct answers of the start and the judge run.}
\label{tab:grpo-pass8-distinct}
\begin{threeparttable}
\begin{tabular}{llcccc}
\toprule
 & & \multicolumn{2}{c}{pass@8} & \multicolumn{2}{c}{Distinct correct} \\
\cmidrule(lr){3-4} \cmidrule(lr){5-6}
Model & Puzzles & Start & Judge run & Start & Judge run \\
\midrule
\multirow{2}{*}{Llama-3.2-3B}
 & 3-number & 0.980 & 0.980 & 2.57 & 2.21 \\
 & 4-number & 0.580 & 0.510 & 1.47 & 1.30 \\
\midrule
\multirow{2}{*}{Qwen3-0.6B}
 & 3-number & 0.880 & 0.755 & 2.08 & 1.60 \\
 & 4-number & 0.320 & 0.275 & 0.62 & 0.48 \\
\bottomrule
\end{tabular}
\begin{tablenotes}\footnotesize
\item The runs of Table~\ref{tab:format-vs-corrupt}: 200 test puzzles per size $\times$ 8 samples at temperature 0.9; start = the SFT model, judge run = format + label-free judge. pass@8: share of puzzles with at least one correct sample. Distinct correct: distinct correct answers per puzzle (whitespace removed).
\end{tablenotes}
\end{threeparttable}
\end{table}


\newpage
\section{Prompts}
\label{app:prompts}
% Generated by data/prompts/make_app_prompts.py from code/countdown/prompts.py (do not edit by hand: rerun it).
\subsection*{Policy prompt}
\noindent\textbf{User}
\begin{Verbatim}[frame=single, fontsize=\small, breaklines, breakanywhere=false]
Using the numbers 25, 4, 7, 3, create an expression that equals 96. You can use +, -, * and /, and each number must be used exactly once.

Respond in the following format:
<reasoning>
{your search for a way to reach the target}
</reasoning>
<solution>
{the steps, one per line, e.g. 25 * 4 = 100}
</solution>
<answer>
{one expression that uses each number exactly once and equals the target}
</answer>
\end{Verbatim}

\subsection*{Judge prompt}
\noindent\textbf{System}
\begin{Verbatim}[frame=single, fontsize=\small, breaklines, breakanywhere=false]
Please solve the following problem step-by-step, making sure that each step is well-justified.
\end{Verbatim}
\noindent\textbf{User}
\begin{Verbatim}[frame=single, fontsize=\small, breaklines, breakanywhere=false]
Using the numbers 25, 4, 7, 3, create an expression that equals 96. You can use +, -, * and /, and each number must be used exactly once.
\end{Verbatim}
\noindent\textbf{Assistant (scored)}
\begin{Verbatim}[frame=single, fontsize=\small, breaklines, breakanywhere=false]
25 * 4 = 100
100 - 7 = 93
93 + 3 = 96
Answer: 25 * 4 - 7 + 3
\end{Verbatim}

\subsection*{Solution from corrupted reasoning}
The policy prompt above, then the assistant turn starts with:
\noindent\textbf{Assistant (forced start)}
\begin{Verbatim}[frame=single, fontsize=\small, breaklines, breakanywhere=false]
<reasoning>
{kept part of the reasoning, empty for no reasoning}
</reasoning>
<solution>
\end{Verbatim}


\noindent Chat templates: Llama-3.2-3B's template starts every prompt with a system turn holding ``Cutting Knowledge Date: December 2023'' and ``Today Date: 26 Jul 2024'' (the date pinned), followed by the system prompt when there is one. Qwen3-0.6B runs with thinking off: with \texttt{enable\_thinking=False}, its template starts the assistant turn with an empty thinking block, \texttt{<think>} and \texttt{</think>} with a blank line between them.

\end{appendices}


\end{document}